% ============================================================================
%  Chapitre 5, partie IV : Estimation par moindres carrés
% ============================================================================

% ----------------------------------------------------------------------------
\section{Estimation par moindres carrés}
% ----------------------------------------------------------------------------

\subsection{Modèle}

Pour la mesure $k$, on empile les deux équations de la partie
précédente :
\[
    \mat{\Phi}_k\,\vect{\theta} = \vect{y}_k,
    \qquad
    \mat{\Phi}_k = \begin{pmatrix}
        \vect{\varphi}_1(\vect{\gamma}_k, \vect{\omega}_k)^T\\
        \vect{\varphi}_2(\vect{\gamma}_k, \vect{\omega}_k)^T
    \end{pmatrix} \in \R^{2\times 8},
    \qquad
    \vect{y}_k = \begin{pmatrix} 0\\ y_2(\vect{\gamma}_k, \vect{\omega}_k)\end{pmatrix}
\]
Sur les $n$ mesures, on empile à nouveau les $\mat{\Phi}_k$ et les
$\vect{y}_k$ en $\mat{\Phi} \in \R^{2n\times 8}$ et
$\vect{y} \in \R^{2n}$ :
\[
    \mat{\Phi}\,\vect{\theta} = \vect{y}
\]

\subsection{Fonction de perte}

On cherche donc le $\vect{\theta}$ qui minimise l'erreur quadratique
sur toutes les mesures
\[
    \mathcal{L}(\vect{\theta})
    = \sum_{k=1}^{n}\frac12\|\vect{y}_k - \mat{\Phi}_k\vect{\theta}\|^2
    = \frac12\|\vect{y} - \mat{\Phi}\vect{\theta}\|^2
\]
Son gradient vaut
\[
    \nabla\mathcal{L}(\vect{\theta})
    = -\mat{\Phi}^T(\vect{y} - \mat{\Phi}\vect{\theta})
\]
et on obtient l'estimation de nos paramètres par : 
\[
    \mat{\Phi}^T\mat{\Phi}\,\vect{\theta} = \mat{\Phi}^T\vect{y}
    \quad\Longrightarrow\quad
    \hat{\vect{\theta}} = \bigl(\mat{\Phi}^T\mat{\Phi}\bigr)^{-1}\mat{\Phi}^T\vect{y}
\]

\subsection{Retour aux moments principaux}

De $\hat{\vect{\theta}}$, on reconstruit $\hat{\mat{J}}$, symétrique,
puis on la diagonalise comme dans la partie I :
\[
    \hat{\mat{J}} = \hat{\mat{P}}\,\mathrm{diag}(\hat{I}_1, \hat{I}_2, \hat{I}_3)\,\hat{\mat{P}}^T
\]
